\documentclass[11pt,a4paper]{article}

% ─────────────────────────── Packages ───────────────────────────
\usepackage[utf8]{inputenc}
\usepackage[T1]{fontenc}
\usepackage[french]{babel}
\usepackage[a4paper,margin=2.4cm]{geometry}
\usepackage{graphicx}
\usepackage{xcolor}
\usepackage{float}
\usepackage{enumitem}
\usepackage{booktabs}
\usepackage{caption}
\usepackage{fancyhdr}
\usepackage[hidelinks]{hyperref}

% ─────────────────────────── Style ───────────────────────────
\definecolor{primary}{HTML}{1F3A5F}
\definecolor{accent}{HTML}{2A7ADE}
\definecolor{light}{HTML}{F2F5FA}

\captionsetup{font=small,labelfont={bf,color=primary}}
% Titres de section colorés (sans package additionnel)
\makeatletter
\renewcommand{\section}{\@startsection{section}{1}{\z@}{-3.5ex plus -1ex minus -.2ex}{2.3ex plus .2ex}{\Large\bfseries\color{primary}}}
\renewcommand{\subsection}{\@startsection{subsection}{2}{\z@}{-3.25ex plus -1ex minus -.2ex}{1.5ex plus .2ex}{\large\bfseries\color{accent}}}
\makeatother

\pagestyle{fancy}
\fancyhf{}
\rhead{\small\color{primary}Analyse de Données de Vente}
\lhead{\small\color{primary}Big Data}
\cfoot{\thepage}
\renewcommand{\headrulewidth}{0.4pt}

% Commande pour insérer une capture centrée
\newcommand{\capture}[3][0.9]{%
  \begin{figure}[H]\centering
  \includegraphics[width=#1\textwidth]{image/#2}
  \caption{#3}\end{figure}}

% ─────────────────────────── Document ───────────────────────────
\begin{document}

% ===== Page de garde =====
\begin{titlepage}
\centering
\vspace*{1.5cm}
{\color{accent}\rule{\textwidth}{2pt}}\\[0.6cm]
{\Huge\bfseries\color{primary} Projet Big Data\par}
\vspace{0.4cm}
{\LARGE Analyse de Données de Vente\par}
\vspace{0.2cm}
{\large\itshape Apache Spark (PySpark) \& Hadoop HDFS sous Docker\par}
{\color{accent}\rule{\textwidth}{2pt}}\\[1.2cm]

{\large\textbf{Dataset :} Online Retail (UCI) — $\approx$ 541\,909 transactions\par}
\vspace{2.2cm}

{\large\bfseries\color{primary} Membres du groupe\par}
\vspace{0.4cm}
\begin{tabular}{l}
\large Alaa Belhadj\\[4pt]
\large Zackaria Ajli\\[4pt]
\large Oussema Jebaari\\[4pt]
\large Aziz Ben Guirat\\[4pt]
\large Mohamed Hamdaoui\\[4pt]
\large Ibtihal El Bannouri\\
\end{tabular}

\vfill
{\large Année universitaire 2025--2026\par}
\end{titlepage}

\tableofcontents
\newpage

% ===== 1. Introduction =====
\section{Introduction et contexte}
Une entreprise de commerce en ligne souhaite exploiter l'historique de ses ventes
afin de mieux comprendre son activité, d'analyser le comportement de ses clients
et d'identifier des opportunités d'amélioration de ses performances commerciales.

L'objectif de ce projet est de développer une \textbf{solution d'analyse de données}
reposant sur \textbf{Apache Spark} pour traiter efficacement plusieurs centaines de
milliers de transactions, puis de stocker le jeu de données sur un cluster
\textbf{Hadoop HDFS}. L'ensemble est conteneurisé avec \textbf{Docker Compose}.

\subsection{Jeu de données}
Le projet s'appuie sur le \textit{Online Retail Dataset} de l'UCI Machine Learning
Repository\footnote{\url{https://archive.ics.uci.edu/dataset/352/online+retail}},
qui contient environ \textbf{541\,909} transactions d'un site de vente en ligne
britannique entre décembre 2010 et décembre 2011. Chaque ligne décrit un article
acheté : numéro de facture, code et description du produit, quantité, date, prix
unitaire, identifiant client et pays.

\begin{table}[H]\centering\small
\begin{tabular}{@{}lll@{}}
\toprule
\textbf{Colonne} & \textbf{Type} & \textbf{Description}\\
\midrule
InvoiceNo   & string  & Numéro de facture\\
StockCode   & string  & Code produit\\
Description & string  & Libellé du produit\\
Quantity    & int     & Quantité achetée\\
InvoiceDate & date    & Date et heure de la transaction\\
UnitPrice   & double  & Prix unitaire (£)\\
CustomerID  & int     & Identifiant client\\
Country     & string  & Pays du client\\
\bottomrule
\end{tabular}
\caption{Schéma du jeu de données Online Retail.}
\end{table}

% ===== 2. Architecture =====
\section{Architecture de la solution}
La solution est composée de trois services orchestrés par \texttt{docker-compose.yml}
sur un réseau commun \texttt{bigdata-net} :

\begin{table}[H]\centering\small
\begin{tabular}{@{}llll@{}}
\toprule
\textbf{Service} & \textbf{Image Docker} & \textbf{Rôle} & \textbf{Interface}\\
\midrule
namenode & bde2020/hadoop-namenode & Maître HDFS (métadonnées) & localhost:9870\\
datanode & bde2020/hadoop-datanode & Stockage des blocs HDFS & localhost:9864\\
spark    & jupyter/pyspark-notebook & Spark + Jupyter (PySpark) & localhost:8888\\
\bottomrule
\end{tabular}
\caption{Les trois conteneurs de la stack Big Data.}
\end{table}

Le dossier local \texttt{data/} (contenant le CSV) est monté dans le conteneur Spark,
et le conteneur Spark écrit ses résultats sur le cluster HDFS via l'adresse
\texttt{hdfs://namenode:9000}.

% ===== 3. Mise en route =====
\section{Mise en route : de \texttt{docker compose up} à l'exécution}

\subsection{Préparation des données}
Le dataset est fourni au format Excel (\texttt{.xlsx}). Un script Python le convertit
en CSV, format directement exploitable par Spark.
\capture[0.75]{01_convert_csv.png}{Conversion du fichier \texttt{.xlsx} en \texttt{.csv}.}

\subsection{Démarrage de la stack}
La commande \texttt{docker compose up -d} construit le réseau, les volumes persistants
et démarre les trois conteneurs.
\capture[0.75]{02_compose_up.png}{Démarrage des conteneurs avec \texttt{docker compose up}.}

On vérifie ensuite que les trois services sont sains (\textit{healthy}).
\capture[1.0]{03_compose_ps.png}{État des conteneurs : \texttt{docker compose ps}.}

\subsection{Accès à Jupyter}
On récupère l'URL d'accès à Jupyter (avec son token de sécurité).
\capture[1.0]{04_jupyter_url.png}{Récupération de l'URL Jupyter.}

% ===== 4. Partie 1 : Spark =====
\section{Partie 1 — Application d'analyse avec Spark}

\subsection{Initialisation de la session Spark}
On crée une \texttt{SparkSession}, point d'entrée de toute application PySpark.
\capture[0.85]{06_code_init.png}{Code — création de la session Spark.}
\capture[0.6]{07_out_init.png}{Sortie — Spark 3.5.0 démarré.}

\subsection{Chargement des données}
Le fichier CSV est chargé dans un \texttt{DataFrame} Spark avec inférence du schéma.
\capture[0.75]{08_code_load.png}{Code — chargement du CSV.}
\capture[0.65]{09_out_load.png}{Sortie — 541\,909 lignes chargées et schéma inféré.}

\subsection{Nettoyage et transformation}
Le jeu de données brut contient des valeurs manquantes (notamment 135\,080 lignes
sans \texttt{CustomerID}) et des lignes correspondant à des retours (quantités
négatives). On applique donc plusieurs traitements :
\begin{itemize}[noitemsep]
  \item suppression des lignes sans identifiant client ;
  \item suppression des quantités et prix négatifs ou nuls ;
  \item suppression des doublons ;
  \item ajout d'une colonne \texttt{TotalPrice} $=$ \texttt{Quantity} $\times$ \texttt{UnitPrice} ;
  \item conversion de la date et extraction de l'année et du mois.
\end{itemize}
\capture[0.9]{11_out_nulls.png}{Sortie — comptage des valeurs nulles par colonne.}
\capture[0.95]{10_code_clean.png}{Code — nettoyage et transformation.}
\capture[1.0]{12_out_clean.png}{Sortie — 392\,692 lignes après nettoyage.}

\subsection{Analyses des ventes et des produits}

\paragraph{Chiffre d'affaires total.}
\capture[0.55]{13_out_ca_total.png}{Chiffre d'affaires total : 8\,887\,208.89 £.}

\paragraph{Produits les plus vendus.}
\capture[0.8]{14_out_top_produits.png}{Top 10 des produits par quantité vendue.}
\capture[0.8]{15_out_top_ca.png}{Top 10 des produits par chiffre d'affaires.}

\paragraph{Répartition géographique.}
Le marché est très largement dominé par le Royaume-Uni (\textgreater 80\,\% du CA).
\capture[0.6]{16_out_pays.png}{CA par pays (Top 10).}
\capture[0.85]{23_chart_pays.png}{Visualisation du CA par pays (hors Royaume-Uni).}

\paragraph{Évolution temporelle.}
On observe une forte croissance du CA à l'approche de la période de Noël 2011.
\capture[0.6]{17_out_mensuel.png}{Évolution mensuelle du CA.}
\capture[0.85]{22_chart_mensuel.png}{Visualisation de l'évolution mensuelle du CA.}

\paragraph{Meilleurs clients.}
\capture[0.55]{18_out_clients.png}{Top 10 des clients par chiffre d'affaires.}

\subsection{Indicateurs clés (KPIs)}
\capture[0.7]{19_out_kpis.png}{Indicateurs synthétiques de l'activité.}

\begin{table}[H]\centering\small
\begin{tabular}{@{}lr@{}}
\toprule
\textbf{Indicateur} & \textbf{Valeur}\\
\midrule
Chiffre d'affaires total      & 8\,887\,208.89 £\\
Commandes uniques             & 18\,532\\
Clients uniques               & 4\,338\\
Produits uniques              & 3\,665\\
Panier moyen par commande     & 479.56 £\\
Articles moyens par commande  & 278.0\\
\bottomrule
\end{tabular}
\caption{Synthèse des indicateurs calculés.}
\end{table}

% ===== 5. Partie 2 : HDFS =====
\section{Partie 2 — Stockage sur Hadoop HDFS}
Conformément à l'énoncé, le jeu de données est stocké sur le cluster HDFS. Deux
versions sont écrites depuis Spark :
\begin{itemize}[noitemsep]
  \item le \textbf{dataset brut} dans \texttt{/data/raw/online\_retail} (CSV) ;
  \item les \textbf{données nettoyées} dans \texttt{/data/online\_retail\_clean}
        (format colonne Parquet, compressé Snappy, 8 partitions).
\end{itemize}
\capture[0.8]{20_code_hdfs.png}{Code — écriture et relecture sur HDFS.}
\capture[1.0]{21_out_hdfs.png}{Sortie — confirmation des écritures HDFS.}

On vérifie enfin l'arborescence stockée sur le cluster :
\capture[1.0]{05_hdfs_ls.png}{Contenu du cluster HDFS (\texttt{hdfs dfs -ls -R /data}).}

La réplication est fixée à 1 (un seul DataNode), et l'interface web du NameNode
(\url{http://localhost:9870}) permet de superviser le cluster et de naviguer dans
le système de fichiers.
\capture[1.0]{24_hdfs_webui.jpeg}{Interface web du NameNode Hadoop (\texttt{localhost:9870}) :
vue d'ensemble du cluster HDFS (version 3.2.1, \textit{Safemode off}, capacité et
utilisation du système de fichiers).}

% ===== 6. Conclusion =====
\section{Conclusion}
Ce projet a permis de mettre en place une chaîne complète de traitement Big Data :
ingestion d'un dataset de plus de 540\,000 transactions, nettoyage et
transformation distribués avec \textbf{PySpark}, calcul d'analyses et d'indicateurs
métier pertinents, puis stockage durable sur un cluster \textbf{Hadoop HDFS}, le tout
reproductible grâce à \textbf{Docker Compose}.

Les analyses montrent une activité fortement concentrée sur le Royaume-Uni, une
saisonnalité marquée (pic à l'approche de Noël) et une dépendance à un petit nombre
de clients à forte valeur. Comme \textbf{extension}, la solution pourrait être étendue
au traitement en temps réel (\textit{streaming}) à l'aide de Spark Structured
Streaming alimenté par Apache Kafka.

\end{document}
